% GENERATED FILE. Edit canonical lessons, then run npm run generate:books.
% canonical-source-hash: fnv1a64:4c88af1741c26e45
% canonical-lessons: ES-C758-voleibol, ES-C758-boxeo, ES-C758-esquiar, ES-C758-patinar, ES-C758-aficionado

\chapter{Más deportes}
\label{ch:es-mas-deportes}
\hlchaptermodality{758}

\begin{chapteropening}
\textbf{By the end of this chapter:} \emph{I can talk about volleyball, boxing, skiing and skating, and about being a fan.}

Complete the last lesson of chapter 758: I can talk about volleyball, boxing, skiing and skating, and about being a fan.
\end{chapteropening}

\section[el voleibol]{el voleibol — volleyball}

\label{lesson:ES-C758-voleibol}



\begin{quote}
Before the new one: say the Spanish for defeat, then the Spanish for the medal.
\end{quote}

\subsection*{You'll want to know: el voleibol}
\textbf{el voleibol} — \textquotedblleft{}volleyball\textquotedblright{}.

\begin{quote}
Jugamos al voleibol en la playa.
\end{quote}

\begin{quote}
\emph{We play volleyball on the beach.}
\end{quote}

Also \emph{el vóley}. In some countries \emph{el vóleibol}.

\begin{cousinweb}
\emph{voleibol} was borrowed from English \emph{volleyball} (EtymDB).
\end{cousinweb}

\begin{grammarlens}[title={What you've built}]
One more word for telling, explaining and describing at B1.
\end{grammarlens}

\subsection*{Your turn}
\begin{itemize}
\raggedright
  \item \emph{Say it:} \emph{el voleibol}
  \item \emph{Say it:} \emph{el voleibol}, once more
  \item \emph{Say it:} a sentence of your own with \emph{el voleibol}
  \item \emph{Recall:} say the Spanish for defeat, then the Spanish for the medal, then say \emph{el voleibol} again
\end{itemize}

\begin{tcolorbox}[breakable,colback=teal!4,colframe=teal!35!black,title={Before you move on}]
What does el voleibol mean? (\textquotedblleft{}Volleyball\textquotedblright{}.) Say it once more.
\end{tcolorbox}
\section[el boxeo]{el boxeo — boxing}

\label{lesson:ES-C758-boxeo}



\begin{quote}
Before the new one: say the Spanish for the medal, then the Spanish for volleyball.
\end{quote}

\subsection*{You'll want to know: el boxeo}
\textbf{el boxeo} — \textquotedblleft{}boxing\textquotedblright{}.

\begin{quote}
No me gusta ver el boxeo.
\end{quote}

\begin{quote}
\emph{I do not like watching boxing.}
\end{quote}

The boxer is \emph{el boxeador}.

\begin{cousinweb}
\emph{boxeo} is built on \emph{boxear}, which was borrowed from English \emph{box} \textquotedblleft{}to fight with the fists\textquotedblright{} (EtymDB).
\end{cousinweb}

\begin{grammarlens}[title={What you've built}]
One more word for telling, explaining and describing at B1.
\end{grammarlens}

\subsection*{Your turn}
\begin{itemize}
\raggedright
  \item \emph{Say it:} \emph{el boxeo}
  \item \emph{Say it:} \emph{el boxeo}, once more
  \item \emph{Say it:} a sentence of your own with \emph{el boxeo}
  \item \emph{Recall:} say the Spanish for the medal, then the Spanish for volleyball, then say \emph{el boxeo} again
\end{itemize}

\begin{tcolorbox}[breakable,colback=teal!4,colframe=teal!35!black,title={Before you move on}]
What does el boxeo mean? (\textquotedblleft{}Boxing\textquotedblright{}.) Say it once more.
\end{tcolorbox}
\section[esquiar]{esquiar — to ski}

\label{lesson:ES-C758-esquiar}



\begin{quote}
Before the new one: say the Spanish for volleyball, then the Spanish for boxing.
\end{quote}

\subsection*{You'll want to know: esquiar}
\textbf{esquiar} — \textquotedblleft{}to ski\textquotedblright{}.

\begin{quote}
En invierno vamos a esquiar.
\end{quote}

\begin{quote}
\emph{In winter we go skiing.}
\end{quote}

Forms: \emph{esquío}, \emph{esquías}. \emph{El esquí} — skiing, a ski.

\begin{cousinweb}
\emph{esquiar} is built on \emph{esquí}, which is usually traced to Norwegian \emph{ski} \textquotedblleft{}a strip of wood\textquotedblright{}, partly through French; Spanish added \emph{e-} before \emph{s}.
\end{cousinweb}

\begin{grammarlens}[title={What you've built}]
One more word for telling, explaining and describing at B1.
\end{grammarlens}

\subsection*{Your turn}
\begin{itemize}
\raggedright
  \item \emph{Say it:} \emph{esquiar}
  \item \emph{Say it:} \emph{esquiar}, once more
  \item \emph{Say it:} a sentence of your own with \emph{esquiar}
  \item \emph{Recall:} say the Spanish for volleyball, then the Spanish for boxing, then say \emph{esquiar} again
\end{itemize}

\begin{tcolorbox}[breakable,colback=teal!4,colframe=teal!35!black,title={Before you move on}]
What does esquiar mean? (\textquotedblleft{}To ski\textquotedblright{}.) Say it once more.
\end{tcolorbox}
\section[patinar]{patinar — to skate; to skid}

\label{lesson:ES-C758-patinar}



\begin{quote}
Before the new one: say the Spanish for boxing, then the Spanish for to ski.
\end{quote}

\subsection*{You'll want to know: patinar}
\textbf{patinar} — \textquotedblleft{}to skate; to skid\textquotedblright{}.

\begin{quote}
Los niños patinan en el parque.
\end{quote}

\begin{quote}
\emph{The children skate in the park.}
\end{quote}

A car can also \emph{patinar} on ice: to skid.

\begin{cousinweb}
\emph{patinar} is built on \emph{patín} \textquotedblleft{}skate\textquotedblright{} (EtymDB), which was borrowed from French \emph{patin}, built on \emph{patte} \textquotedblleft{}paw, foot\textquotedblright{}.
\end{cousinweb}

\begin{grammarlens}[title={What you've built}]
One more word for telling, explaining and describing at B1.
\end{grammarlens}

\subsection*{Your turn}
\begin{itemize}
\raggedright
  \item \emph{Say it:} \emph{patinar}
  \item \emph{Say it:} \emph{patinar}, once more
  \item \emph{Say it:} a sentence of your own with \emph{patinar}
  \item \emph{Recall:} say the Spanish for boxing, then the Spanish for to ski, then say \emph{patinar} again
\end{itemize}

\begin{tcolorbox}[breakable,colback=teal!4,colframe=teal!35!black,title={Before you move on}]
What does patinar mean? (\textquotedblleft{}To skate; to skid\textquotedblright{}.) Say it once more.
\end{tcolorbox}
\section[el aficionado]{el aficionado — the fan; the amateur}

\label{lesson:ES-C758-aficionado}



\begin{quote}
Before the new one: say the Spanish for to ski, then the Spanish for to skate.
\end{quote}

\subsection*{You'll want to know: el aficionado}
\textbf{el aficionado} — \textquotedblleft{}the fan; the amateur\textquotedblright{}.

\begin{quote}
Los aficionados cantan en el estadio.
\end{quote}

\begin{quote}
\emph{The fans sing in the stadium.}
\end{quote}

Also an amateur: \emph{un fotógrafo aficionado}.

\begin{cousinweb}
\emph{aficionado} is the participle of \emph{aficionar} (EtymDB), built on \emph{afición} \textquotedblleft{}liking\textquotedblright{}, which comes from Latin \emph{affectiō}.
\end{cousinweb}

\begin{grammarlens}[title={What you've built}]
One more word for telling, explaining and describing at B1.
\end{grammarlens}

\subsection*{Your turn}
\begin{itemize}
\raggedright
  \item \emph{Say it:} \emph{el aficionado}
  \item \emph{Say it:} \emph{el aficionado}, once more
  \item \emph{Say it:} a sentence of your own with \emph{el aficionado}
  \item \emph{Recall:} say the Spanish for to ski, then the Spanish for to skate, then say \emph{el aficionado} again
\end{itemize}

\begin{tcolorbox}[breakable,colback=teal!4,colframe=teal!35!black,title={Before you move on}]
What does el aficionado mean? (\textquotedblleft{}The fan; the amateur\textquotedblright{}.) Say it once more.
\end{tcolorbox}
